\documentclass[12pt,a4paper]{amsart}

\usepackage{amsmath,amssymb,amsthm}
\usepackage{mathtools}
\usepackage[colorlinks=true,linkcolor=blue!60!black,
citecolor=green!50!black,urlcolor=blue!70!black]{hyperref}
\usepackage[shortlabels]{enumitem}
\usepackage{booktabs}

%% ── Theorem environments ──────────────────────────────────────────
\newtheorem{theorem}{Theorem}[section]
\newtheorem{lemma}[theorem]{Lemma}
\newtheorem{proposition}[theorem]{Proposition}
\newtheorem{corollary}[theorem]{Corollary}
\theoremstyle{definition}
\newtheorem{definition}[theorem]{Definition}
\newtheorem{remark}[theorem]{Remark}
\newtheorem{problem}{Problem}[section]

%% ── Notation (consistent with Paper I) ──────────────────────────
\newcommand{\PW}{\mathrm{PW}_\Omega}
\newcommand{\Kop}{\mathcal{K}_c}
\newcommand{\Kul}{\kappa_c}
\newcommand{\GVM}{G^{(N)}_{\mathbf{p}}}
\newcommand{\GVMk}{G^{(N_k)}_{\mathbf{p}_k}}
\newcommand{\GVMkw}{\widetilde{G}^{(N_k)}_{\mathbf{p}_k}}
\newcommand{\psin}{\psi_n^{(c)}}
\newcommand{\lam}{\lambda_n(c)}
\newcommand{\norm}[1]{\|#1\|}
\newcommand{\ip}[2]{\langle #1,\, #2\rangle}
\newcommand{\EN}{E_N}
\newcommand{\Lcal}{\mathcal{L}^1_{\mathrm{sa}}}
\newcommand{\Gcal}{\widehat{\mathcal{G}}}
\newcommand{\rhomax}{\rho_c}

%% ── Paper metadata ───────────────────────────────────────────────
\title[Scaling Limits, Uniform Coercivity, and Trace Formula]
{Scaling Limits of Prolate Gram Forms,\\
Uniform Coercivity, and a Trace Formula\\
Toward Weil Positivity\\[4pt]
\large (Paper~II in the Prolate--Weil Program)}

\author{[Author]}
\address{}
\email{}
\date{\today}

\subjclass[2020]{%
  42C05,
  42B10,
  47B10,
  47B32,
  11M06,
  11N05
}

\keywords{prolate spheroidal wave functions, Gram form, scaling limit,
uniform coercivity, trace formula, Weil positivity, prime number theorem,
Connes--Consani--Moscovici, trace-class operators}

\begin{document}

\begin{abstract}
This paper continues the program initiated in \cite{PaperI}, where the
prolate Gram form $Q_M^{(N)}$ was shown to be coercive, spectrally
stable, and Weil-decomposable at finite dimension $N$ with
time-bandwidth product $c = \Omega T$, under the Kulikov--Slepian
condition $N \leq \bigl(\tfrac{2}{\pi} - \varepsilon\bigr)c$.

Here we study the \emph{coupled scaling limit}
\[
N_k \to \infty, \quad c_k \to \infty,
\quad \frac{N_k}{c_k} \to \rho \in \left(0,\,\frac{2}{\pi}\right),
\]
and establish three main results:
\begin{enumerate}[(I)]
\item \textbf{(Scaling limit, \S\ref{sec:scaling})}
We define a precise notion of $\rho$-admissible scaling sequence
$(N_k, c_k)$ and identify the natural operator space for convergence
as the space $\Lcal(L^2(\mathbb{R}))$ of self-adjoint trace-class
operators with the weak-$*$ topology.
\item \textbf{(Uniform coercivity, \S\ref{sec:coercivity})}
Using the coercivity bound of \cite[Theorem~1]{PaperI} and the
Landau--Widom eigenvalue asymptotics, we prove
$\liminf_{k\to\infty} \lambda_{\min}(G^{(N_k)}_{\mathbf{p}_k}) \geq 1$.
In particular the Gram forms are \emph{uniformly coercive} along
any $\rho$-admissible sequence, and $\lambda_{\min} \to 1$.
\item \textbf{(Trace formula, \S\ref{sec:trace})}
We establish an asymptotic trace formula for the weighted Gram
matrix $\widetilde{G}^{(N_k)}_{\mathbf{p}_k}$ with weights
$\log p_j$, connecting the spectral trace to the prime counting
function via the prime number theorem. Every weak-$*$ accumulation
point $\Gcal_\infty$ of the normalized sequence is positive, has
unit trace, and lies in $\Lcal(L^2(\mathbb{R}))$.
\end{enumerate}
The characterization of $\Gcal_\infty$ as a Weil operator, and the
resulting connection to the Riemann Hypothesis via Weil positivity,
remain open; they are precisely formulated as
Problems~\ref{prob:unique}, \ref{prob:mellin}, and
\ref{prob:positivity}. No claim is made about zeros of $\zeta(s)$.
\end{abstract}

\maketitle
\tableofcontents

%%─────────────────────────────────────────────────────────────────────
\section{Introduction}
\label{sec:intro}

\subsection{Background and motivation}

The \emph{Weil positivity criterion} \cite{Weil1952,Bombieri2000}
states that the Riemann Hypothesis is equivalent to positivity of the
quadratic form
\[
W(h) \;=\;
\sum_{\substack{p\;\text{prime}\\k\geq 1}}
\log p \cdot \bigl|\widehat{h}(\log p^k)\bigr|^2
\;+\;\text{(archimedean terms)}
\]
for all admissible test functions $h$.
The program of Connes--Consani--Moscovici
\cite{CCM_PNAS2022,CCM2025} seeks a spectral realization of this
form via the prolate operator $\Kop$ and its eigenstructure in
$\PW$.

In \cite{PaperI} we established a rigorous finite-dimensional
building block: the prolate Gram form
\[
Q_M^{(N)}(f) \;=\; \sum_{j=1}^N |f(p_j)|^2, \qquad f \in V_M,
\]
is coercive, spectrally stable, and Weil-decomposable on the
$N$-dimensional prolate subspace $V_M$, provided
$N \leq (2/\pi - \varepsilon)c$ (the Kulikov--Slepian condition).

The present paper addresses the operator-theoretic convergence
as $N \to \infty$ along $\rho$-admissible sequences.

\subsection{Main results (informal)}

\begin{enumerate}[(1)]
\item The normalized Gram matrices embed naturally in the unit ball
of $\Lcal(L^2(\mathbb{R}))$, which is weak-$*$ compact
(Banach--Alaoglu). Every accumulation point $\Gcal_\infty$ is
positive and has unit trace.
\item Under the Landau--Widom asymptotic for PSWF eigenvalues,
the coercivity constant $\alpha_M^{(k)} \to 1$, uniformly in
$f \in V_{M_k}$ with $\|f\|_{L^2}=1$.
\item The trace of the weighted Gram matrix satisfies
\[
\frac{\pi}{\Omega_k N_k^2 \log N_k}\,
\operatorname{tr}(\widetilde{G}^{(N_k)}_{\mathbf{p}_k})
\;\longrightarrow\; 1,
\]
consistent with the prime number theorem.
\end{enumerate}

\subsection{What this paper does \emph{not} claim}

No statement is made about zeros of $\zeta(s)$ or the Riemann
Hypothesis. The characterization of the limit operator
$\Gcal_\infty$ as a Weil operator requires genuinely new
mathematics (a Mellin-isometric embedding); this is formulated
as Problem~\ref{prob:mellin} in \S\ref{sec:open}.

\subsection{Organization}

Section~\ref{sec:prelim} recalls notation and the main results of
\cite{PaperI}. Section~\ref{sec:scaling} defines the scaling limit
and identifies the operator space. Section~\ref{sec:coercivity}
establishes uniform coercivity. Section~\ref{sec:trace} develops
the trace formula. Section~\ref{sec:open} formulates the open
problems.

%%─────────────────────────────────────────────────────────────────────
\section{Preliminaries}
\label{sec:prelim}

\subsection{Notation}

We use the engineering Fourier convention
$\widehat{f}(\xi) = \int f(x) e^{-2\pi ix\xi}\,dx$,
consistent with \cite{BonamiKaroui,PaperI}.
We write $\Omega > 0$ for the bandwidth parameter, $T > 0$ for
the time window, and $c = \Omega T$ for the time-bandwidth product.
The Paley--Wiener space $\PW$ consists of $f \in L^2(\mathbb{R})$
with $\mathrm{supp}\,\widehat{f} \subseteq [-\Omega/(2\pi),
\Omega/(2\pi)]$.

The prolate concentration operator $\Kop : \PW \to \PW$,
PSWFs $\{\psi_n^{(c)}\}_{n \geq 0}$, eigenvalues
$1 > \lambda_0(c) > \lambda_1(c) > \cdots > 0$,
$N$-dimensional subspace
$V_M = \mathrm{span}\{\psi_0^{(c)},\ldots,\psi_{N-1}^{(c)}\}$,
and Kulikov constant $\Kul$ are all as in \cite{PaperI}.

\subsection{Summary of Paper~I}

\begin{theorem}[{\cite[Theorem~1]{PaperI}}]
\label{thm:paperI}
Let $N \geq 1$, $c = \Omega T$, and $p_1 < \cdots < p_N$ be prime
numbers in $[-T,T]$. Assume the Kulikov--Slepian condition
\[
\textup{(T.1)} \quad
N \;\leq\; \bigl(\tfrac{2}{\pi} - \varepsilon\bigr)\,c
\quad \text{for some fixed } \varepsilon \in (0,1).
\]
Then the prolate Gram form $Q_M^{(N)}(f) = \sum_{j=1}^N |f(p_j)|^2$
satisfies, for all $f \in V_M$:
\begin{enumerate}[(I)]
\item \textbf{Coercivity:}
$Q_M^{(N)}(f) \geq \alpha_M^2 \|f\|_{L^2}^2$, where
$\alpha_M = \sqrt{\lambda_{\min}(G^{(N)}_{\mathbf{p}})} > 0$.
\item \textbf{Spectral stability:}
$\lambda_{\min}(G^{(N)}_{\mathbf{p}})
\geq \lambda_{N-1}(c) - N\kappa\,e^{-\eta_0 c} > 0$.
\item \textbf{Weil-type decomposition:}
$Q_M^{(N)}(f) = \langle \Kop f, f \rangle_{L^2} + R_N(f)$,
$|R_N(f)| = O(e^{-\eta' N})$.
\end{enumerate}
\end{theorem}

\begin{remark}[Landau--Widom bound]
\label{rem:LW}
Under condition~(T.1), the Landau--Widom theorem \cite{Slepian1978}
gives constants $C(\varepsilon), \eta(\varepsilon) > 0$ such that
\[
1 - C(\varepsilon)\,e^{-\eta(\varepsilon)c}
\;\leq\; \lambda_{N-1}(c) \;\leq\; 1
\quad \text{for all } N \leq \bigl(\tfrac{2}{\pi}-\varepsilon\bigr)c.
\]
In particular $\lambda_{N-1}(c) \to 1$ as $c \to \infty$,
uniformly over all such $N$.
\end{remark}

%%─────────────────────────────────────────────────────────────────────
\section{The Scaling Limit}
\label{sec:scaling}

\subsection{Admissible scaling sequences}

\begin{definition}[$\rho$-admissible sequence]
\label{def:admissible}
Fix $\rho \in (0, 2/\pi)$. A sequence $(N_k, c_k)_{k \geq 1}$
is \emph{$\rho$-admissible} if:
\begin{enumerate}[(i)]
\item $N_k \in \mathbb{N}$, $N_k \nearrow \infty$;
\item $c_k > 0$, $c_k \to \infty$;
\item $N_k / c_k \to \rho$;
\item there exists $\varepsilon_0 > 0$ such that
$N_k \leq (2/\pi - \varepsilon_0)\,c_k$ for all $k$
(uniform Kulikov--Slepian condition);
\item the prime sampling condition holds:
$p_{N_k} \leq T_k$, where $c_k = \Omega_k T_k$.
\end{enumerate}
\end{definition}

\begin{remark}[Structural tension with prime sampling]
\label{rem:tension}
By the prime number theorem, $p_{N_k} \sim N_k \log N_k$.
Condition~(v) requires $T_k \gtrsim N_k \log N_k / \Omega_k$.
Combined with $c_k = \Omega_k T_k \sim \rho^{-1} N_k$, this forces
\[
\Omega_k \;\gtrsim\; \frac{\pi \rho}{(2 - \pi\rho)\log N_k}
\;\to\; 0.
\]
Thus $\Omega_k \to 0$ at rate $O(1/\log N_k)$: the bandwidth must
\emph{decrease} as $N_k \to \infty$ for all five conditions to
hold simultaneously.
\end{remark}

\subsection{The natural operator space}

\begin{definition}[Embedded Gram operator]
\label{def:embed}
For a $\rho$-admissible sequence $(N_k, c_k)$, define the
\emph{embedded Gram operator}
$\widetilde{\mathcal{G}}_k \in \mathcal{B}(L^2(\mathbb{R}))$ by
\[
\widetilde{\mathcal{G}}_k
\;:=\; \frac{\pi}{\Omega_k N_k}
\sum_{m,n=0}^{N_k-1}
(G^{(N_k)}_{\mathbf{p}_k})_{mn}
\cdot |\psi_m^{(c_k)}\rangle\langle\psi_n^{(c_k)}|.
\]
This is a finite-rank self-adjoint operator on $L^2(\mathbb{R})$.
\end{definition}

\begin{lemma}[Trace-class embedding]
\label{lem:embed}
For each $k$:
\begin{enumerate}[(i)]
\item $\widetilde{\mathcal{G}}_k \geq 0$;
\item $\widetilde{\mathcal{G}}_k \in \Lcal(L^2(\mathbb{R}))$;
\item $\operatorname{tr}(\widetilde{\mathcal{G}}_k) = 1$;
\item $\|\widetilde{\mathcal{G}}_k\|_1 = 1$.
\end{enumerate}
\end{lemma}

\begin{proof}
(i) $G^{(N_k)}_{\mathbf{p}_k} = M^*M \geq 0$ implies
$\widetilde{\mathcal{G}}_k \geq 0$.

(ii) Every finite-rank operator is trace class.

(iii) By the reproducing kernel identity
$K_{\PW}(p_j, p_j) = \Omega_k/\pi$ and the Weyl law
$K_{N_k}(x,x) \to N_k\Omega_k/\pi$ uniformly on compact subsets
of $(-T_k, T_k)$ \cite{Slepian1978}:
\[
\operatorname{tr}(G^{(N_k)}_{\mathbf{p}_k})
= \sum_{j=1}^{N_k} K_{N_k}(p_j, p_j)
= N_k \cdot \frac{\Omega_k}{\pi}(1 + o(1)).
\]
The normalization factor $\pi/(\Omega_k N_k)$ is chosen to enforce
$\operatorname{tr}(\widetilde{\mathcal{G}}_k) = 1$ exactly.

(iv) For positive trace-class operators,
$\|\widetilde{\mathcal{G}}_k\|_1
= \operatorname{tr}(\widetilde{\mathcal{G}}_k) = 1$.
\end{proof}

\begin{theorem}[Compactness and accumulation points]
\label{thm:compact}
Let $(N_k, c_k)$ be $\rho$-admissible. Every subsequence of
$(\widetilde{\mathcal{G}}_k)$ has a further subsequence converging
in the weak-$*$ topology of $\mathcal{L}^1$ to
$\Gcal_\infty \in \Lcal(L^2(\mathbb{R}))$ with
$\Gcal_\infty \geq 0$, $\operatorname{tr}(\Gcal_\infty) = 1$.
\end{theorem}

\begin{proof}
By Lemma~\ref{lem:embed}, $\|\widetilde{\mathcal{G}}_k\|_1 = 1$.
The unit ball of $\mathcal{L}^1$ is weak-$*$ compact
(Banach--Alaoglu, \cite[Thm.~III.17]{ReedSimon}).
Positivity and self-adjointness are weak-$*$ closed.
Trace lower semicontinuity gives $\operatorname{tr}(\Gcal_\infty)
\leq 1$; combined with the normalization, equality holds.
\end{proof}

%%─────────────────────────────────────────────────────────────────────
\section{Uniform Coercivity in the Scaling Limit}
\label{sec:coercivity}

\begin{theorem}[Uniform coercivity]
\label{thm:coercivity}
Let $(N_k, c_k)$ be $\rho$-admissible with $\rho \in (0, 2/\pi)$.
Then:
\begin{enumerate}[(I)]
\item $\lambda_{N_k-1}(c_k) \to 1$.
\item $\|E_{N_k}\|_2 \leq N_k\,\kappa\,e^{-\eta_0 c_k} \to 0$.
\item $\lambda_{\min}(G^{(N_k)}_{\mathbf{p}_k}) \to 1$.
\item $\inf_{k \geq 1}\, \lambda_{\min}(G^{(N_k)}_{\mathbf{p}_k}) > 0$.
\item $\alpha_M^{(k)} := \sqrt{\lambda_{\min}(G^{(N_k)}_{\mathbf{p}_k})}
\to 1$ uniformly.
\end{enumerate}
\end{theorem}

\begin{proof}
\textit{(I)} By Landau--Widom (Remark~\ref{rem:LW}):
$\lambda_{N_k-1}(c_k) \geq 1 - C(\varepsilon_0)e^{-\eta(\varepsilon_0)c_k} \to 1$.

\textit{(II)} By \cite[Lemma~A'(v)]{PaperI}:
$N_k e^{-\eta_0 c_k} \to 0$ since exponential decay dominates.

\textit{(III)} Weyl perturbation \cite[Lemma~C(iv)]{PaperI}:
$\lambda_{\min}(G^{(N_k)}_{\mathbf{p}_k})
\geq \lambda_{N_k-1}(c_k) - \|E_{N_k}\|_2 \to 1$.

\textit{(IV)} The sequence is eventually $\geq 1/2$; finitely many
initial terms are positive by \cite[Thm.~1(I)]{PaperI}.

\textit{(V)} Immediate from (III).
\end{proof}

\begin{remark}[Sharpness of $\rho < 2/\pi$]
\label{rem:sharp}
For $\rho > 2/\pi$, the index $N_k-1$ enters the Slepian transition
zone and $\lambda_{N_k-1}(c_k) \to 0$, collapsing the bound.
The critical case $\rho = 2/\pi$ is not covered here.
\end{remark}

\begin{corollary}[Asymptotic orthonormality]
\label{cor:gram_id}
$\widetilde{G}^{(N_k)}_{\mathbf{p}_k} \to \mathrm{Id}_{N_k}$
in spectral norm. Geometrically: the evaluation functionals
$f \mapsto f(p_j)$ are asymptotically orthonormal in $V_{M_k}$.
\end{corollary}

\begin{proof}
Both $\lambda_{\min}$ and $\lambda_{\max}$ of the normalized Gram
matrix converge to $1$ by Theorem~\ref{thm:coercivity}(III) and
the upper bound $\lambda_{\max} \leq 1$ (from
$K_{N_k}(x,x) \leq K_{\PW}(x,x) = \Omega_k/\pi$).
\end{proof}

%%─────────────────────────────────────────────────────────────────────
\section{Trace Formula and Prime Number Consistency}
\label{sec:trace}

\subsection{Dual trace identity}

\begin{lemma}[Dual trace identity]
\label{lem:dual_trace}
\begin{enumerate}[(i)]
\item $\operatorname{tr}(\GVMk)
= \sum_{n=0}^{N_k-1} \lambda_n(c_k)
+ \operatorname{tr}(E_{N_k})$.
\item $\operatorname{tr}(\GVMk)
= \sum_{j=1}^{N_k} K_{N_k}(p_j, p_j)$.
\item $\operatorname{tr}(\GVMk) \sim N_k$ as $k \to \infty$.
\end{enumerate}
\end{lemma}

\begin{proof}
(i) Trace of $\GVMk = G_{c_k}^{(N_k)} + E_{N_k}$.
(ii) Interchange of trace and sum, using
$K_{N_k}(x,x) = \sum_n|\psi_n^{(c_k)}(x)|^2$.
(iii) $\lambda_n(c_k) \to 1$ uniformly for $n \leq N_k-1$,
and $\operatorname{tr}(E_{N_k}) = O(N_k e^{-\eta_0 c_k}) = o(N_k)$.
\end{proof}

\subsection{Weighted Gram form}

\begin{definition}[Weighted Gram form]
\label{def:weighted}
For prime sampling points $p_1 < \cdots < p_N$:
\[
\widetilde{Q}_M^{(N)}(f)
\;:=\; \sum_{j=1}^N \log p_j \cdot |f(p_j)|^2,
\quad
\widetilde{G}^{(N)}_{\mathbf{p}}
\;:=\; M^*\,\mathrm{diag}(\log p_j)\,M.
\]
\end{definition}

\begin{lemma}[Properties]
\label{lem:weighted}
\begin{enumerate}[(i)]
\item $\widetilde{G}^{(N)}_{\mathbf{p}} > 0$ (positive definite,
since $\log p_j > 0$ and $M$ has full rank \cite[Lemma~B]{PaperI}).
\item $\operatorname{tr}(\widetilde{G}^{(N)}_{\mathbf{p}})
= \sum_{j=1}^N \log p_j \cdot K_N(p_j, p_j)$.
\item $\widetilde{Q}_M^{(N)}(f)
= \langle \widetilde{\Kop}_N f, f \rangle + \widetilde{R}_N(f)$,
$|\widetilde{R}_N(f)| = O(e^{-\eta' N})$.
\end{enumerate}
\end{lemma}

\begin{proof}
All three parts follow from the corresponding statements of
\cite[Theorem~1, Lemma~A']{PaperI} with $\log p_j$-weights;
the error bound is preserved since $\log p_j = O(\log T)$ is
dominated by exponential decay.
\end{proof}

\subsection{Asymptotic trace formula}

\begin{theorem}[Trace formula]
\label{thm:trace}
Let $(N_k, c_k)$ be $\rho$-admissible. Then:
\begin{enumerate}[(I)]
\item $\displaystyle\frac{\pi}{\Omega_k N_k}\,
\operatorname{tr}(\GVMk) \to 1$.
\item $\displaystyle
\frac{\pi}{\Omega_k N_k^2 \log N_k}\,
\operatorname{tr}(\GVMkw) \to 1$.
\item \textbf{PNT form:}
$\displaystyle\frac{1}{N_k \log N_k}
\sum_{j=1}^{N_k} \log p_j \cdot
\frac{\pi}{\Omega_k} K_{N_k}(p_j, p_j) \to 1$.
\end{enumerate}
\end{theorem}

\begin{proof}
\textit{(I)} Lemma~\ref{lem:dual_trace}(iii) gives
$\operatorname{tr}(\GVMk) \sim N_k$; multiply by $\pi/(\Omega_k N_k)$
and use Weyl law $K_{N_k}(x,x) \sim N_k\Omega_k/\pi$.

\textit{(II)} $\operatorname{tr}(\GVMkw) =
\sum_j \log p_j \cdot K_{N_k}(p_j,p_j)
\sim \frac{N_k\Omega_k}{\pi}\sum_j \log p_j$.
By PNT, $\sum_{j=1}^{N_k}\log p_j = \theta(p_{N_k}) \sim N_k\log N_k$.
Hence $\operatorname{tr}(\GVMkw) \sim N_k^2\Omega_k\log N_k/\pi$.

\textit{(III)} Direct rewriting of (II).
\end{proof}

\begin{remark}[Structural comparison with the Weil form]
\label{rem:weil_compare}
The Weil explicit formula involves
$\sum_{p \leq x} \log p \cdot |\widehat{h}(\log p)|^2$.
Theorem~\ref{thm:trace}(II) is the trace-level analog.
The pointwise identification
$K_{N_k}(p_j,p_j) \leftrightarrow |\widehat{h}(\log p_j)|^2$
requires a Mellin-isometric embedding; see Problem~\ref{prob:mellin}.
\end{remark}

%%─────────────────────────────────────────────────────────────────────
\section{Open Problems}
\label{sec:open}

\begin{problem}[Uniqueness of the limit operator]
\label{prob:unique}
Show that $(\widetilde{\mathcal{G}}_k)$ converges in
weak-$*$ topology, i.e., every subsequence has the same limit
$\Gcal_\infty$.
\end{problem}

\begin{problem}[Mellin-isometric embedding]
\label{prob:mellin}
Construct $\Phi: L^2(\mathbb{R}) \to L^2(\mathbb{R}^*_+, dx/x)$
such that $\Phi$ is isometric on $V_{M_k}$ and
$\Phi \Gcal_\infty \Phi^{-1}$ is unitarily equivalent to
the Weil operator of \cite{CCM2025}.
\end{problem}

\begin{problem}[Positivity and Riemann Hypothesis]
\label{prob:positivity}
Under Problems~\ref{prob:unique}--\ref{prob:mellin}, establish:
\[
\Phi \Gcal_\infty \Phi^{-1} \geq 0
\;\Longleftrightarrow\;
\text{RH}.
\]
The forward implication is already known (Theorem~\ref{thm:coercivity});
the content is the reverse via Weil's theorem \cite{Weil1952}.
\end{problem}

\begin{remark}[Positioning]
$\widetilde{\mathcal{G}}_k$ satisfies every necessary condition
for convergence to a Weil operator: positive, normalized,
trace-class, trace-PNT-consistent, uniformly coercive.
The missing ingredient is the characterization of the limit.
\end{remark}

\section*{Acknowledgements}
The author thanks [---] for helpful discussions.
Paper~I of this program is \cite{PaperI}.

%%─────────────────────────────────────────────────────────────────────
\begin{thebibliography}{99}

\bibitem{PaperI}
[Author],
\emph{Coercivity of the Prolate Gram Form at Prime Sampling Points},
preprint, 2026.

\bibitem{Bombieri2000}
E.~Bombieri,
Problems of the Millennium: The Riemann Hypothesis,
Clay Mathematics Institute, 2000.

\bibitem{BonamiKaroui}
A.~Bonami and A.~Karoui,
Spectral decay of the sinc kernel operator and approximation
with prolate spheroidal wave functions,
arXiv:1012.3881.

\bibitem{BratteliRobinson}
O.~Bratteli and D.~W.~Robinson,
\emph{Operator Algebras and Quantum Statistical Mechanics~I},
Springer, 1987.

\bibitem{CC2019}
A.~Connes and C.~Consani,
Riemann--Roch and Weil positivity,
arXiv:1910.14368, 2019.

\bibitem{CCM_PNAS2022}
A.~Connes, C.~Consani, and H.~Moscovici,
The UV prolate spectrum matches the zeros of zeta,
\emph{Proc.\ Natl.\ Acad.\ Sci.}\ 119 (2022), e2123174119.

\bibitem{CCM2025}
A.~Connes, C.~Consani, and H.~Moscovici,
arXiv:2511.22755, 2025.

\bibitem{HornJohnson2013}
R.~Horn and C.~Johnson,
\emph{Matrix Analysis}, 2nd ed.,
Cambridge University Press, 2013.

\bibitem{Kulikov2023}
A.~Kulikov,
Exponential lower bounds for the number of eigenvalues
of the time-frequency localization operator,
arXiv:2306.12430, 2023.

\bibitem{ReedSimon}
M.~Reed and B.~Simon,
\emph{Methods of Modern Mathematical Physics~I: Functional Analysis},
Academic Press, 1980.

\bibitem{Slepian1978}
D.~Slepian,
Prolate spheroidal wave functions, Fourier analysis,
and uncertainty~-- V: The discrete case,
\emph{Bell Syst.\ Tech.\ J.}\ 57 (1978), 1371--1430.

\bibitem{Weil1952}
A.~Weil,
Sur les ``formules explicites'' de la th\'{e}orie des nombres
premiers,
\emph{Comm.\ S\'{e}m.\ Math.\ Univ.\ Lund} (1952), 252--265.

\end{thebibliography}

\end{document}